\section*{Abstract}

$\xipar$ is the single fundamental parameter of T0 theory. The complete chain reads: $\xipar \to \{m_e, m_\mu, \ldots\} \to G \to \lP$, with $L_0 = \xipar \cdot \lP$ as sub-Planck scale of spacetime granulation. The maximal scale $R_H$ follows from $E_H = \Ezero \cdot \xipar^{41/4}$. The ratio $R_H/L_0 = 6{.}49 \times 10^{64}$ follows entirely from $\xipar = 4/30000$ alone.

The question of a universal maximum size of the universe is, however, not meaningfully posed in FFGFT: \textbf{every physical scale possesses its own system-specific upper bound} $R_\text{torus}(m) = \hbar/mc$. There is no single cosmic upper bound valid for all systems — $R_H$ is the upper bound of the cosmological sector, not of the universe as a whole.

\section{Derivation Chain: From \texorpdfstring{$\xi$}{xi} to Both Scale Limits}

\subsection{The Single Fundamental Parameter}

T0 theory has one single fundamental parameter (Document 013):
\begin{equation}
	\boxed{\xipar = \frac{4}{3} \times 10^{-4}}
\end{equation}
All physical quantities derive from it. The complete chain reads:
\begin{equation}
	\xipar \;\to\; \{m_e, m_\mu, m_\tau, \ldots\} \;\to\; G \;\to\; \lP \;\to\; c
\end{equation}
with $L_0 = \xipar \cdot \lP$ as the fundamental sub-Planck scale of spacetime granulation.

\subsection{Derivation of Particle Masses from \texorpdfstring{$\xi$}{xi}}

The lepton masses follow from $\xipar$ via the geometric quantum number factor $f(n,l,j)$ and the T0 scaling factor $S_{T0} = 1\,\text{MeV}/c^2$:
\begin{equation}
	m_\ell = \frac{f(n,l,j)^2}{\xipar^2} \cdot S_{T0}
\end{equation}
The geometric mean of the lepton masses yields the characteristic energy scale:
\begin{equation}
	\Ezero = \sqrt{m_e \cdot m_\mu} = 7{.}398\,\text{MeV}
\end{equation}
from which the fine-structure constant follows without free parameters:
\begin{equation}
	\alpha = \xipar \cdot \frac{\Ezero^2}{(1\,\text{MeV})^2} = \frac{1}{137{.}036} \quad \checkmark
\end{equation}

\subsection{Gravitational Constant and Planck Length}

The gravitational constant follows from $\xipar$ and $m_e$ in natural units (Document 013):
\begin{equation}
	G = \frac{\xipar^2}{4\,m_e} \qquad\Rightarrow\qquad \lP = \sqrt{G} = \frac{\xipar}{2\sqrt{m_e}}
\end{equation}
In natural units with $m_e = 0{.}511\,\text{MeV}$: $\lP \approx 9{.}33 \times 10^{-5}$ (natural units), which after SI conversion gives:
\begin{equation}
	\lP = 1{.}616 \times 10^{-35}\,\text{m}
\end{equation}

\subsection{Minimal Length Scale}

\begin{keyresult}[Minimal Length Scale from \texorpdfstring{$\xi$}{xi}]
\begin{equation}
	\boxed{L_0 = \xipar \cdot \lP = 2{.}155 \times 10^{-39}\,\text{m}}
\end{equation}
The complete derivation chain $\xipar \to m_e \to G \to \lP \to L_0$ depends only on $\xipar$ as the single fundamental input. SI numerical values arise through the unavoidable calibration to human measurement units — physically the theory is parameter-free.
\end{keyresult}

\section{The Maximal Scale: Hubble Radius from \texorpdfstring{$\xi$}{xi}}

\subsection{Hubble Energy from \texorpdfstring{$\xi$}{xi}}

The Hubble energy follows directly from $\xipar$ (Document 026):
\begin{equation}
	E_H = \Ezero \cdot \xipar^{41/4}
\end{equation}
with $\Ezero = \sqrt{m_e \cdot m_\mu} = 7{.}398\,\text{MeV}$ and $\xipar = 4/30000$.

Numerical evaluation:
\begin{align}
	\xipar^{41/4} &= (1{.}333 \times 10^{-4})^{41/4} = 1{.}908 \times 10^{-40} \\
	E_H &= 7{.}398\,\text{MeV} \times 1{.}908 \times 10^{-40} = 1{.}412 \times 10^{-33}\,\text{eV}
\end{align}

\begin{keyresult}[Hubble Constant from \texorpdfstring{$\xi$}{xi}]
\begin{equation}
	\boxed{H_0^{\text{T0}} = \frac{E_H}{\hbar} = \frac{1{.}412 \times 10^{-33}\,\text{eV}}{6{.}582 \times 10^{-16}\,\text{eV\,s}} \approx 66{.}2\,\text{km/s/Mpc}}
\end{equation}
Deviation from the Planck value ($67{.}4\,\text{km/s/Mpc}$): $-1{.}9\%$.
\end{keyresult}

\subsection{Hubble Radius as Maximal Scale}

\begin{foundation}[Distinction: Hubble radius vs.\ particle horizon]
Standard cosmology ($\Lambda$CDM, expanding universe) distinguishes two cosmic horizons: the \textit{Hubble radius} $R_H = c/H_0 \approx 1{.}4 \times 10^{26}$\,m and the substantially larger \textit{particle horizon} (observable universe, $\approx 4{.}4 \times 10^{26}$\,m $\approx 46.5$\,Gly), which encompasses the full causal past since the Big Bang. The particle horizon is only meaningfully defined in an expanding universe with a Big Bang.

In T0 theory there is no expansion and no Big Bang. The Hubble radius $R_H$ is the sole relevant geometric horizon: beyond $R_H$ the accumulated fractal path elongation yields $z \to \infty$, i.e.\ no information can be received. A particle horizon in the $\Lambda$CDM sense does not exist conceptually in T0. Numerical comparisons with standard values must therefore always refer to $R_H$, not to the particle horizon.
\end{foundation}

The Hubble radius follows directly from the Hubble energy:
\begin{equation}
	R_H = \frac{\hbar c}{E_H} = \frac{1{.}973 \times 10^{-7}\,\text{eV\,m}}{1{.}412 \times 10^{-33}\,\text{eV}}
\end{equation}

\begin{keyresult}[Hubble Radius from \texorpdfstring{$\xi$}{xi}]
\begin{equation}
	\boxed{R_H = 1{.}398 \times 10^{26}\,\text{m}}
\end{equation}
Derived entirely from $\xipar$, without free parameters.
\end{keyresult}

\subsection{Hubble Radius and Total Mass}

The total mass of the observable universe corresponding to the Hubble radius:
\begin{equation}
	M_U = \frac{R_H \cdot c^2}{2G} = 9{.}43 \times 10^{52}\,\text{kg} = 4{.}7 \times 10^{22}\,M_\odot
\end{equation}

\begin{keyresult}[Hubble Radius from \texorpdfstring{$\xi$}{xi}]
\begin{equation}
	\boxed{R_H = 1{.}398 \times 10^{26}\,\text{m}}
\end{equation}
Derived entirely from $\xipar$, without free parameters.
\end{keyresult}

\section{The Complete Scale Hierarchy}

\subsection{Scale Hierarchy}

\begin{table}[h]
\centering
\begin{tabular}{lll}
\toprule
\textbf{Scale} & \textbf{Value} & \textbf{Interpretation} \\
\midrule
$L_0 = \xipar \cdot \lP$ & $2{.}155 \times 10^{-39}\,\text{m}$ & Sub-Planck scale \\
$\lP$ & $1{.}616 \times 10^{-35}\,\text{m}$ & Planck length \\
$\lambda_e = \hbar/m_e c$ & $3{.}86 \times 10^{-13}\,\text{m}$ & Electron Compton wavelength \\
$R_H = \hbar c / E_H$ & $1{.}398 \times 10^{26}\,\text{m}$ & Hubble horizon \\
\bottomrule
\end{tabular}
\caption{Complete scale hierarchy from $\xipar = 4/30000$}
\end{table}

\subsection{The Ratio of Scales}

\begin{equation}
	\frac{R_H}{L_0} = \frac{1{.}398 \times 10^{26}\,\text{m}}{2{.}155 \times 10^{-39}\,\text{m}}
\end{equation}

\begin{keyresult}[Scale Ratio from \texorpdfstring{$\xi$}{xi}]
\begin{equation}
	\boxed{\frac{R_H}{L_0} = 6{.}49 \times 10^{64}}
\end{equation}
57 orders of magnitude between the minimal and maximal scale — both from $\xipar = 4/30000$.
\end{keyresult}

\subsection{The Complete Derivation Chain}

\begin{keyresult}[Complete Chain from \texorpdfstring{$\xi$}{xi} to Both Limits]
\begin{align}
	&\text{Lagrangian} \;\Rightarrow\; \xipar = 4/30000 \notag \\
	&\quad\Downarrow \notag \\
	&G_{\text{SI}} = \frac{\xipar^2}{4m_e} \times C_{\text{dim}} \times C_{\text{conv}} \times \Kfrak \;\Rightarrow\; \lP = \sqrt{\hbar G/c^3} \notag \\
	&\quad\Downarrow\qquad\qquad\qquad\quad\Downarrow \notag \\
	&L_0 = \xipar \cdot \lP \qquad\qquad\qquad E_H = \Ezero \cdot \xipar^{41/4} \notag \\
	&= 2{.}155 \times 10^{-39}\,\text{m} \qquad\quad\Downarrow \notag \\
	&\qquad\qquad\qquad R_H = \hbar c/E_H \notag \\
	&\qquad\qquad\qquad = 1{.}398 \times 10^{26}\,\text{m} \notag \\
	&\quad\Downarrow \notag \\
	&\boxed{R_H/L_0 = 6{.}49 \times 10^{64} \quad \text{(from } \xipar \text{ alone)}}
\end{align}
\end{keyresult}

\section{Physical Meaning}

\subsection{Physical Meaning of \texorpdfstring{$R_H$}{R\_H}}

In a static universe without expansion the Hubble radius is the geometric horizon beyond which photons in the $\xipar$ field are completely redshifted ($z \to \infty$) through accumulated fractal path elongation — not through intrinsic energy loss of the photon itself, but through geometric stretching of the propagation path. Tired light does not occur.

\begin{keypoint}[Redshift without Expansion and without Frequency Dependence]
The redshift arises through fractal summation of the path extension of photons in the T0 vacuum field: $z \approx \xipar \cdot \ln(d/\lP)$. Since the path extension is a property of the geometry itself — not of the photon — all frequencies experience exactly the same stretching. There is no frequency dependence. The effect is observationally equivalent to an expanding universe but mechanistically fundamentally different. Derived in detail in Document 026 (Geometric Cosmology).
\end{keypoint}

\subsection{Scale relativity of the torus: fractal nesting}

In T0 theory the 4D torus is not a cosmic container in which particles reside. Rather, every physical system carries \textbf{its own characteristic torus}, whose scale is set by the Compton wavelength $\bar\lambda = \hbar/mc$ of that system:

\begin{equation}
	R_{\text{torus}}(m) = \frac{\hbar}{mc}
\end{equation}

\begin{table}[h]
\centering
\begin{tabular}{lll}
\toprule
\textbf{System} & \textbf{Characteristic scale} & \textbf{Ratio to $L_0$} \\
\midrule
Electron     & $3{.}86 \times 10^{-13}$\,m & $\sim 10^{26}$ \\
Proton       & $2{.}10 \times 10^{-16}$\,m & $\sim 10^{23}$ \\
Hydrogen     & $5{.}29 \times 10^{-11}$\,m & $\sim 10^{28}$ \\
Galaxy       & $\sim 10^{21}$\,m           & $\sim 10^{60}$ \\
Universe     & $R_H = 1{.}40 \times 10^{26}$\,m & $6{.}49 \times 10^{64}$ \\
\bottomrule
\end{tabular}
\caption{System-specific torus scales — all sharing $L_0$ as common lower bound}
\end{table}

\begin{keypoint}[Fractal nesting without rigid hierarchy]
These structures are not nested like Russian dolls — the electron torus does not sit \textit{inside} the galactic torus. Each scale is the geometric self-structure of the respective system, constituted by its mass. The common lower bound $L_0$ is the same at every level; the upper bound is system-specific.

The fractal coupling between scales is mediated by $\xipar$ and therefore very small. Interactions between widely separated scales — for instance between electron structure and the cosmological field — are not zero, but suppressed by powers of $\xipar \approx 1{.}33 \times 10^{-4}$. This weak cross-scale coupling is physically relevant: it is the origin of the fractal correction factors $K_{\text{frak}}$ in the calculation of $g$-factors, particle masses, and $G$ — small but measurable contributions that are absent in the Standard Model or must be introduced as free parameters.
\end{keypoint}

\subsection{Locality, flat torus, and finiteness of calculations}

The T0 torus is not the intuitive donut torus embedded in $\mathbb{R}^3$. It is the \textbf{flat identification torus} $T^n = \mathbb{R}^n / \mathbb{Z}^n$: a Euclidean space with opposite boundaries identified. This torus has \textbf{intrinsic curvature $K = 0$ everywhere pointwise} — not merely in the integral. Locally it is exactly identical to $\mathbb{R}^n$, without any approximation.

\begin{keypoint}[Flat torus: locally exactly Euclidean]
For the embedded donut torus in $\mathbb{R}^3$, Gaussian curvature varies: positive on the outside ($K > 0$), negative on the inside ($K < 0$). The integral vanishes exactly (Gauss-Bonnet: $\int K\,dA = 2\pi\chi = 0$ for $\chi = 0$). For the flat identification torus, by contrast, $K = 0$ holds \emph{everywhere} — there is no inner and outer curvature that must cancel in an integral. Calculations within this torus are exactly the same as in unbounded Euclidean space. No curvature corrections, no topological terms.
\end{keypoint}

Three consequences follow for calculations in T0 theory:

\textbf{1. No renormalisation required.} All field integrals run over the finite local torus with scale $R_\text{torus}(m) = \hbar/mc$. Since this torus is finite and locally flat, UV divergences are structurally excluded. The upper integration limit is not $\infty$ but $R_\text{torus}$; the lower limit is $L_0$. Both are finite and derived from $\xipar$.

\textbf{2. The $\xipar$ correction for an infinite universe is not negligibly small — it is simply irrelevant.} Formally an infinite universe would generate a correction $\sim \xipar \cdot \ln(R_\infty / R_H)$, which diverges logarithmically as $R_\infty \to \infty$. This divergence does not apply, however, because all physical observables are integrals over the \emph{local} torus of the respective system. Contributions from scales $\gg R_\text{torus}(m)$ do not couple to local physics — they are suppressed by powers of $\xipar$ and lie outside the integration domain. Whether the universe is finite or infinite has no effect on local calculations.

\textbf{3. Curvature corrections are absolutely negligible.} The relative curvature correction at the smallest physical scale $L_0$ within a torus of scale $R_\text{torus}$ is $(L_0 / R_\text{torus})^2$:
\begin{align}
	\text{Electron:} \quad & \left(\frac{L_0}{\bar\lambda_e}\right)^2 = \left(\frac{2{.}155\times 10^{-39}}{3{.}86\times 10^{-13}}\right)^2 \approx 3{.}1 \times 10^{-53} \notag \\
	\text{Proton:} \quad & \left(\frac{L_0}{\bar\lambda_p}\right)^2 \approx 1{.}1 \times 10^{-46} \notag
\end{align}
These values are exactly zero for all practical calculations.

\begin{keypoint}[One field — masses as knots — emergent coherence length]
Behind these three consequences lies a unified ontological structure: there is \textbf{one single $\xipar$ vacuum field}. There is no separate hierarchy of fields for different scales.

Masses are \textbf{topological knots} in this field — localised structures, not external objects. Each knot of mass $m$ carries, through the duality $\tilde{T} = 1/m$, a characteristic time scale, and from this its spatial coherence length $\hbar/mc$ emerges immediately. The torus radius is not externally assigned but follows from the knot structure itself.

The hierarchy of torus scales — from the proton knot ($\sim 10^{-16}$\,m) to the cosmological horizon ($\sim 10^{26}$\,m) — is therefore nothing other than the \textbf{fractal self-similarity of a single field}: the same $\xipar$ at every scale, the same lower bound $L_0$, system-specific coherence lengths as emergent properties of the respective knot structure. The cross-scale coupling is not zero, but suppressed by powers of $\xipar$ — which is why one calculates locally without knowing the full field and still obtains exact results.
\end{keypoint}

The exponent $41/4$ in $E_H = \Ezero \cdot \xipar^{41/4}$ follows from the renormalisation group structure of the T0 field over cosmological scales and still requires full derivation from the $\xipar$ field theory (Document 026). The numerical agreement with $H_0 = 66{.}2\,\text{km/s/Mpc}$ ($-1{.}9\%$ from Planck) is a strong consistency argument.

\begin{foundation}[On the role of $R_H$ in calculations]
The precise numerical value of $R_H$ plays a subordinate role in the physical calculations of T0 theory. The structurally significant quantity is $L_0 = \xipar \cdot \lP$ — it follows purely from $\xipar$ and is the true foundation of all scale statements.

$R_H$, by contrast, is different in two respects. First, its numerical value is tied to the measured $H_0$ (or depends on the exponent $41/4$ still to be fully derived). Second, $R_H$ in T0 is not primarily a computed result but a \textbf{boundary by definition}: the scale beyond which $z \to \infty$ holds and no observation is possible. More than $R_H$ is structurally inaccessible within the model — regardless of whether the exponent $41/4$ is rigorously derived or fitted to $H_0$.

The ratio $R_H/L_0 = 6{.}49 \times 10^{64}$ therefore does not express a deep theoretical truth, but the distance between the geometrically derived lower bound and the observation-based horizon. The genuine achievement of the theory lies in $L_0$ — not in $R_H$.
\end{foundation}

\subsection{Cosmological Constant from \texorpdfstring{$\xi$}{xi}}

The dimensionless cosmological constant follows directly:
\begin{equation}
	\Lambda_{\text{dim}} = \left(\frac{\lP}{R_H}\right)^2 = \left(\frac{1{.}616 \times 10^{-35}}{1{.}398 \times 10^{26}}\right)^2 \approx 1{.}34 \times 10^{-122}
\end{equation}
The standard value is $\Lambda_{\text{dim}} \approx 2{.}87 \times 10^{-122}$ (from $\Lambda \approx 1{.}1 \times 10^{-52}\,\text{m}^{-2}$, CODATA). The T0 result is in the same order of magnitude (factor~2.1) — not dark energy, but a geometric consequence of the scale hierarchy from $\xipar$. Full agreement requires the still-pending derivation of the exponent $41/4$ (Document 026).

\section{Summary}

\begin{keyresult}[Complete Scale Hierarchy from \texorpdfstring{$\xipar = 4/30000$}{xi = 4/30000}]
\begin{align}
	\text{Lower limit:} \quad & L_0 = \xipar \cdot \lP = 2{.}155 \times 10^{-39}\,\text{m} \notag \\
	\text{Planck:} \quad & \lP = 1{.}616 \times 10^{-35}\,\text{m} \notag \\
	\text{Electron:} \quad & \lambda_e = \hbar/m_e c = 3{.}86 \times 10^{-13}\,\text{m} \notag \\
	\text{Hubble:} \quad & R_H = \hbar c/E_H = 1{.}398 \times 10^{26}\,\text{m} \notag
\end{align}
\begin{equation}
	\boxed{\text{All scales from } \xipar = 4/30000 — \text{no free parameter}}
\end{equation}
The universe is geometrically self-consistently closed — not by a boundary, but by the topology of the 4D torus with $L_0$ and $R_H$ as natural limits of the same geometry.

\begin{keyresult}[Key conclusion: Is there a maximum size of the universe?]
The question is not meaningfully posed as a global question in FFGFT. Every physical scale — from the proton to the cosmological horizon — possesses its \textbf{own system-specific upper bound} $R_\text{torus}(m) = \hbar/mc$, emergent from the knot structure of the respective mass in the single $\xipar$ field. $R_H$ is the upper bound of the cosmological sector. The universal lower bound $L_0$ applies equally to all scales — a universal upper bound does not exist, because there is no scale-independent physics for which it would be relevant.
\end{keyresult}
\end{keyresult}

\begin{foundation}[Epistemological framing: model boundaries, not existence boundaries]
$L_0$ and $R_H$ are not physical walls and make no claim about the total size or global structure of the universe. They are \textbf{boundaries of describability} within T0 theory: below $L_0$ the spacetime granulation diverges, beyond $R_H$ the redshift diverges ($z \to \infty$). Whether anything exists outside these boundaries is simply not a question the theory addresses.

The decisive advantage over other approaches lies not in this epistemological modesty alone, but in the fact that T0 \textbf{requires no regularisation tricks whatsoever}: no renormalisation procedures to remove UV divergences, no artificial cutoffs, no counterterms, no auxiliary fields introduced to tame infinities. $L_0$ does not arise as a post-hoc correction — it emerges as a direct geometric consequence of $\xipar$. The theory is finite from the outset, not because infinities have been subtracted away, but because the parameter $\xipar$ structurally excludes them.
\end{foundation}

\begin{foundation}[Note: Connection to the SI System]
The absolute numerical values in this document — $L_0$, $R_H$, $M_U$ in SI units — presuppose a connection to the SI unit system. This is not circular but epistemologically necessary: from a dimensionless $\xipar$ alone, no absolute metre or kilogram values can be derived. The connection is made via $m_e$, $\lP$ or $\alpha$ as equivalent entry points. Treated in detail in Document 180 (Section: Connection to the SI System) and in Documents 056 and 101.

Documents 180, 181 and 182 form a unit and should not be read separately.
\end{foundation}

\vspace{1cm}
\noindent\textit{The exponent $41/4$ still requires full derivation from the $\xipar$ field theory (Document 026).}
